% Standalone ResearchIQ application architecture figure for the ITLT variant.
% Build: xelatex researchiq_app_architecture.tex
\documentclass[border=6pt]{standalone}
\usepackage{fontspec}
\usepackage{polyglossia}
\setmainlanguage{ukrainian}
\setotherlanguage{english}
\setmainfont{DejaVu Serif}
\setsansfont{DejaVu Sans}
\setmonofont{DejaVu Sans Mono}
\usepackage{tikz}
\usetikzlibrary{arrows.meta,positioning,fit}
\begin{document}
\begin{tikzpicture}[
  font=\normalsize,
  node distance=7mm and 16mm,
  box/.style={draw, rounded corners, align=center, inner sep=5pt, minimum height=11mm},
  wide/.style={box, minimum width=38mm},
  smallbox/.style={box, minimum width=34mm},
  arrow/.style={-{Stealth[length=2.4mm]}, thick},
  dashedarrow/.style={-{Stealth[length=2.4mm]}, thick, dashed}
]
  % Left column: React SPA on top, the manager/admin (block 1) stacked below it
  % (block 2). Stacking narrows the whole figure so the boxes — and their text —
  % render larger when the figure is fit to the text width.
  \node[wide, fill=cyan!8] (spa) {\textbf{React SPA}\\\footnotesize dashboard, НПП, ORCID search\\chat, reports};
  \node[wide, below=of spa, fill=blue!8] (user) {\textbf{Керівник / адміністратор}\\\footnotesize запити, контроль, експорт};
  \node[wide, right=of spa, fill=orange!10] (api) {\textbf{Node.js / Express API}\\\footnotesize auth, tenant routing\\data and report endpoints};

  % API sits at the top-left of the server cluster so the SPA→API arrow enters
  % the container at its boundary instead of crossing the box border.
  \node[smallbox, below=9mm of api, fill=gray!10] (tenant) {\textbf{Tenant storage}\\\footnotesize окремі JSON-набори\\для кожного ЗВО};
  \node[smallbox, right=of tenant, fill=green!10] (er) {\textbf{Golden Record}\\\footnotesize нормалізація, ER\\провенанс, HITL};
  \node[smallbox, right=of er, fill=purple!10] (agent) {\textbf{\texttt{pi-agent-core}}\\\footnotesize tool calling\\пояснення і звіти};

  \node[smallbox, below=9mm of er, fill=yellow!12] (sources) {\textbf{ORCID / OpenAlex / Scopus}\\\footnotesize профілі, праці, метрики};
  \node[smallbox, below=9mm of agent, fill=red!8] (llm) {\textbf{Gemini / OpenRouter}\\\footnotesize інтерпретація запиту\\без критичних рішень};

  \draw[arrow] (user) -- (spa);
  \draw[arrow] (spa) -- (api);
  \draw[arrow] (api) -- (tenant);
  \draw[arrow] (api) -- (er);
  \draw[arrow] (api) -- (agent);
  \draw[arrow] (sources) -- (er);
  \draw[arrow] (er) -- (tenant);
  \draw[arrow] (agent) -- (er);
  \draw[dashedarrow] (agent) -- (llm);
  \draw[dashedarrow] (llm) -- (agent);

  \node[draw, rounded corners, fit=(api)(tenant)(er)(agent), inner sep=7pt,
        label={[font=\bfseries]above:Серверний контур ResearchIQ}] {};
\end{tikzpicture}
\end{document}
